\documentclass[11pt,a4paper]{article}

% ── Packages ──────────────────────────────────────────────────────
\usepackage[utf8]{inputenc}
\usepackage[T1]{fontenc}
\usepackage{amsmath,amssymb,amsthm}
\usepackage{mathtools}
\usepackage{hyperref}
\usepackage{booktabs}
\usepackage[margin=2.5cm]{geometry}

% ── Theorem environments ─────────────────────────────────────────
\newtheorem{theorem}{Theorem}[section]
\newtheorem{proposition}[theorem]{Proposition}
\newtheorem{conjecture}[theorem]{Conjecture}
\newtheorem{corollary}[theorem]{Corollary}
\newtheorem{lemma}[theorem]{Lemma}
\theoremstyle{definition}
\newtheorem{definition}[theorem]{Definition}
\theoremstyle{remark}
\newtheorem{remark}[theorem]{Remark}

% ── Notation ─────────────────────────────────────────────────────
\DeclareMathOperator{\Ai}{Ai}
\DeclareMathOperator{\Bi}{Bi}
\newcommand{\Vq}{V_{\mathrm{quad}}}
\newcommand{\abs}[1]{\lvert #1 \rvert}
\newcommand{\Z}{\mathbb{Z}}
\newcommand{\Q}{\mathbb{Q}}
\newcommand{\R}{\mathbb{R}}
\newcommand{\C}{\mathbb{C}}

% ── Metadata ─────────────────────────────────────────────────────
\title{A non-classical Painlev\'e~V transcendent
from a quadratic polynomial continued fraction:
surface classification and resurgent Stokes data}
\author{papanokechi}
\date{\today}

\begin{document}
\maketitle

% ══════════════════════════════════════════════════════════════════
%  ABSTRACT
% ══════════════════════════════════════════════════════════════════
\begin{abstract}
We study the resurgence and isomonodromic
structure of a constant $V_{\mathrm{quad}}
\approx 1.19737\ldots$ arising as the limit
of a quadratic polynomial continued fraction
(PCF). The convergent denominators satisfy a
Wallis-type recurrence whose continuum limit
yields an exact second-order ODE of the form
$a(x)y'' + a'(x)y' + c(x)y = 0$, with
$a(x) = 3x^2+x+1$ and $c(x) = -x^2$. This
exact-ODE structure forces both finite
singularities to be apparent and explains,
via monodromy, the failure of CM-period
and hypergeometric identification strategies.
The CM discriminant $\mathrm{disc}(a)
= 1 - 12 = -11$ that appears throughout
the analysis is the discriminant of the
Wallis characteristic polynomial, not a
deep CM structure; this explains why
exhaustive PSLQ searches against CM
period bases at discriminant $-11$ return
null by a structural argument independent
of precision.

The isomonodromic deformation of the
underlying connection is governed by the
fifth Painlev\'e equation with parameters
$\alpha=1/6$, $\beta=\gamma=0$,
$\delta=-1/2$ ($\delta\neq0$); its space of
initial conditions is the Sakai surface
$D_5^{(1)}$ with affine Weyl symmetry
$W(A_3^{(1)})$. Since $\delta\neq0$ the
equation does not reduce to Painlev\'e~III,
and the constant is expected to be a
generic, non-classical PV transcendent.
The formal asymptotic series at infinity
is Gevrey-$1$ divergent; its Borel
transform has a branch-point singularity
at $\xi_0\approx2/\sqrt{3}$ (the Stokes gap)
with branch exponent
$\beta_{\mathrm{exp}}=-1/(3\sqrt{3})$,
confirmed to eight significant figures
by Domb--Sykes analysis and
Richardson$^3$ extrapolation.
The Stokes constant $S \approx
0.43770528\ldots$ is computed to eight
digits by the Dingle late-term formula.
An exhaustive search over $245$ variants
of the Jimbo (1982) connection formula
establishes that $S$ is not a simple
combination of Gamma values at the WKB
parameters; its closed-form identification
requires the connection parameter
$\sigma_{\mathrm{conn}}$, which is
transcendental in the accessory parameter
$q=(5+i\sqrt{11})/54$.

The structure $b(x)=a'(x)$ identified
here is a general feature of PCF-derived
ODEs: for any polynomial continued
fraction, the Wallis recurrence produces
an exact ODE whose apparent singularity
positions are the roots of the Wallis
characteristic polynomial and whose CM
discriminant is $\mathrm{disc}(a)$.

\smallskip
\noindent\textit{AI assistance disclosure.}
$\Vq$ was discovered and
characterized through the SIARC
(Self-Iterating Analytic Relay Chain)
multi-model AI relay pipeline.
Manuscript preparation involved Claude
(\texttt{claude-sonnet-4-6}, Anthropic)
and GitHub Copilot (Microsoft). All
mathematical claims were independently
verified by the author.
\end{abstract}

\noindent
\textbf{2020 Mathematics Subject Classification:}
34M55, 34M56, 34E05, 11J70, 33E30.

\noindent
\textbf{Keywords:}
polynomial continued fractions,
resurgence, Stokes constant,
apparent singularity, confluent Heun equation,
Painlev\'e~V, Sakai surface $D_5^{(1)}$,
Borel summation,
arithmetic identification.


% ══════════════════════════════════════════════════════════════════
%  1. INTRODUCTION
% ══════════════════════════════════════════════════════════════════
\section{Introduction}\label{sec:intro}

A \emph{polynomial continued fraction} (PCF)
of degree~$d$ is an expression
\[
  K = b(0) + \underset{n \ge 1}{\mathbf{K}}\,
  \frac{a(n)}{b(n)},
  \qquad a(n), b(n) \in \Z[n],
  \quad \deg a, \deg b \leq d.
\]
These objects have a distinguished history: the
Gauss continued fraction expresses ${}_2F_1$
ratios~\cite{lorentzen2008}, and Ap\'ery's
proof of the irrationality of
$\zeta(3)$~\cite{apery1979} exploits the PCF
with $a(n) = -n^6$ and
$b(n) = 34n^3 + 51n^2 + 27n + 5$.
Recent systematic searches, notably the
Ramanujan Machine project~\cite{raayoni2021},
have generated new conjectural PCF identities
for known constants.

This paper studies the PCF with constant
partial numerators $a(n)=1$ and quadratic
partial denominators $b(n)=3n^2+n+1$:
\begin{equation}\label{eq:vquad-def}
  \Vq = 1 + \underset{n \ge 1}{K}\,
  \frac{1}{3n^2+n+1}
  = 1.19737399068835760244\ldots
\end{equation}
We proved in~\cite{vquad-exclusion} that $\Vq$
is irrational with irrationality measure
$\mu(\Vq)=2$, and that it is not a rational
linear combination of values from 16
special-function families at 2\,050-digit
precision. The present paper explains
\emph{why} these identification searches fail
via a structural analysis of the governing
differential equation.

\subsection{Main results}

Conceptually, the paper has two intertwined
themes: (i)~a structural mechanism that forces
apparent singularities and obstructs CM and
hypergeometric identification, and (ii)~a
detailed resurgent analysis of the associated
Painlev\'e~V ($D_5^{(1)}$) connection problem.

\begin{enumerate}
\item \textbf{(Exact ODE structure,
  \S\ref{sec:vquad})}
  The convergent denominators satisfy a
  three-term recurrence whose continuum
  limit produces the ODE
  \begin{equation}\label{eq:ode-intro}
    a(x)\,y'' + a'(x)\,y' + c(x)\,y = 0,
  \end{equation}
  where $a(x) = 3x^2+x+1$ and $c(x)=-x^2$.
  The identity $b(x) \equiv a'(x)$ is the
  structural key.

\item \textbf{(Apparent singularity exclusion,
  \S\ref{sec:frobenius})}
  Both finite singularities
  $s_{1,2} = (-1 \pm i\sqrt{11})/6$ are
  apparent: indicial exponents $\{0,0\}$,
  trivial monodromy. This rules out
  identification of $\Vq$ via CM periods or
  classical hypergeometric special values
  (Theorem~\ref{thm:exclusion2}).

\item \textbf{(Painlev\'e classification,
  \S\ref{sec:painleve})}
  The isomonodromic deformation is governed
  by the fifth Painlev\'e equation with
  parameters $(\alpha,\beta,\gamma,\delta)
  = (1/6, 0, 0, -1/2)$; its space of initial
  conditions is the Sakai surface $D_5^{(1)}$
  (symmetry $W(A_3^{(1)})$). Since
  $\delta\neq0$ it does not reduce to
  Painlev\'e~III.

\item \textbf{(Resurgence analysis,
  \S\ref{sec:resurgence})}
  The formal asymptotic series at $\infty$ is
  Gevrey-$1$ divergent. Its Borel transform
  has a branch-point singularity at
  $\xi_0 \approx 2/\sqrt{3}$ (the Stokes gap
  $\sigma_+ - \sigma_-$) with branch exponent
  $\beta_{\mathrm{exp}} = -1/(3\sqrt{3})$,
  confirmed to eight significant figures
  (Theorem~\ref{thm:borel}).

\item \textbf{(Stokes constant,
  \S\ref{sec:stokes})}
  The Stokes constant $S \approx 0.43770528$
  is computed to eight digits via
  Dingle late-term extraction but remains
  unidentified in closed form.
\end{enumerate}


% ══════════════════════════════════════════════════════════════════
%  2. THE CONSTANT V_QUAD
% ══════════════════════════════════════════════════════════════════
\section{The constant $\Vq$}\label{sec:vquad}

\subsection{Definition and convergence}

The convergent denominators $Q_n$ of
the PCF~\eqref{eq:vquad-def} satisfy
the three-term recurrence
\begin{equation}\label{eq:recurrence}
  Q_{n+1} = (3n^2+n+1)\,Q_n + Q_{n-1},
  \qquad Q_{-1}=0,\; Q_0=1.
\end{equation}
Since $b(n) = 3n^2+n+1 \geq 5$ for $n \geq 1$,
the \'Sleszyński--Pringsheim theorem guarantees
convergence. The Wronskian identity
$P_nQ_{n-1} - P_{n-1}Q_n = (-1)^{n+1}$ yields
the irrationality of $\Vq$ with optimal
irrationality measure $\mu(\Vq)=2$.

\subsection{The governing ODE}

\paragraph{Formal continuum passage from \eqref{eq:recurrence}.}
Let $Q(x)$ be a smooth interpolant of the convergent
denominator sequence, so that $Q(n)=Q_n$, and Taylor-expand
about $x=n$:
\begin{equation}\label{eq:taylor-symm}
  Q_{n\pm 1}
  \;=\; Q(x) \pm Q'(x) + \tfrac{1}{2}Q''(x)
  + O\!\bigl(Q'''(x)\bigr).
\end{equation}
Substituting into the recurrence~\eqref{eq:recurrence},
and writing $a(x) := 3x^2+x+1$ for the polynomial
that interpolates the recurrence coefficient
$b(n)=3n^2+n+1$, gives the formal balance
\begin{equation}\label{eq:continuum-step}
  2Q(x) + Q''(x) + O(Q''')
  \;=\; a(x)\,Q(x) + Q_{n-1}
  \;=\; a(x)\,Q(x) + Q(x) - Q'(x) + \tfrac{1}{2}Q''(x) + O(Q'''),
\end{equation}
so that at leading order the first-derivative term
$-Q'(x)$ cannot be absorbed into the $Q$, $Q''$ balance
and must be repackaged against the polynomial
$a'(x)=6x+1$ --- the identity $b(x)=a'(x)$ of
\eqref{eq:exact-identity} below. Differentiating the
leading balance and combining via the exact-form
ansatz $\frac{d}{dx}[a(x)\,Q'] + c(x)\,Q = 0$ yields
\begin{equation}\label{eq:ode}
  a(x)\,y'' + b(x)\,y' + c(x)\,y = 0,
\end{equation}
where
\begin{equation}\label{eq:ode-coeffs}
  a(x) = 3x^2+x+1, \quad
  b(x) = 6x+1, \quad
  c(x) = -x^2.
\end{equation}

\begin{remark}[Confluent Heun class]\label{rem:heunc}
The ODE~\eqref{eq:ode} belongs to the confluent Heun class in the
sense of Ronveaux~\cite{ronveaux1995}: it has exactly two finite
regular singular points (at the roots $s_{1,2}=(-1\pm i\sqrt{11})/6$
of $a(x)=3x^2+x+1$) and one irregular singular point of Poincar\'e
rank~$1$ at $x=\infty$ (verified: $\lim_{x\to\infty}q(x)=-1/3\ne 0$
where $q=-x^2/a(x)$). This singularity structure places the
connection problem and the associated isomonodromic deformation
in the confluent Heun sector, making the Painlev\'e~V ($D_5^{(1)}$)
identification of Theorem~\ref{thm:painleve} available via the
Jimbo--Miwa--Ueno framework~\cite{jimbo1981}. Explicit HeunC
parameters in Ronveaux normal form require a coordinate change
sending $s_1\mapsto 0$, $s_2\mapsto 1$ followed by a gauge
transformation $y=e^{\lambda z}\psi$; we do not use these
parameters further and omit their derivation.
\end{remark}

The trailing coefficient $c(x)=-x^2$ is fixed by
matching the large-$n$ WKB scales
$\sigma_\pm = \pm 1/\sqrt{3}$ of~\eqref{eq:wkb}
to the exponential growth rate of the dominant Wallis
solution $Q_n$; we omit the routine verification.
This continuum limit is formal; the ODE captures the
leading asymptotic behavior of $Q_n$ for large $n$.

The decisive algebraic identity is
\begin{equation}\label{eq:exact-identity}
  b(x) = a'(x)
  \quad\text{i.e.}\quad
  6x+1 \equiv \frac{d}{dx}(3x^2+x+1).
\end{equation}
This holds symbolically and has the immediate
consequence that the ODE is \emph{exact}: the
first two terms form a total derivative,
\begin{equation}\label{eq:exact-ode}
  \frac{d}{dx}\bigl[a(x)\,y'\bigr]
  + c(x)\,y = 0.
\end{equation}

\begin{table}[ht]
\centering
\caption{Polynomial data for the
$\Vq$ continued fraction.}
\label{tab:poly}
\small
\begin{tabular}{@{}llll@{}}
\toprule
Object & Polynomial & Degree & Discriminant \\
\midrule
Partial numerator & $a(n) = 1$ & 0 & --- \\
Partial denominator & $b(n) = 3n^2+n+1$ & 2 & $-11$ \\
ODE leading coeff.\ & $a(x) = 3x^2+x+1$ & 2 & $-11$ \\
ODE middle coeff.\ & $b(x) = 6x+1 = a'(x)$ & 1 & --- \\
ODE trailing coeff.\ & $c(x) = -x^2$ & 2 & $0$ \\
\bottomrule
\end{tabular}
\end{table}

\begin{remark}[PCF-derived ODEs and the
  $b(x)=a'(x)$ structure]
\label{rem:pcf-ode-general}
The identity $b(x) = a'(x)$ is not a
coincidence specific to $\Vq$: it is a
general feature of PCF-derived ODEs.
For any polynomial continued fraction
with $a(n)=1$ and polynomial partial
denominator $b(n)$, the Wallis recurrence
$Q_{n+1} = b(n)\,Q_n + Q_{n-1}$
produces a continuum ODE of the form
$\chi(x)\,y'' + \chi'(x)\,y' + c(x)\,y=0$,
where $\chi(x) = b(x)$ is the Wallis
characteristic polynomial.
The exact-ODE structure
$\frac{d}{dx}[\chi(x)\,y'] + c(x)\,y=0$
is therefore automatic, and the finite
singularities of the ODE --- the roots
of $\chi$ --- are always apparent.

In particular, the ``CM discriminant''
$-11$ that appears throughout the analysis
of $\Vq$ is simply
$\mathrm{disc}(\chi) = \mathrm{disc}(3x^2+x+1)
= 1 - 12 = -11$
--- the discriminant of the characteristic
polynomial. For a different quadratic PCF
with $b(n) = \alpha n^2 + \beta n + \gamma$,
the CM discriminant would be
$\beta^2 - 4\alpha\gamma$, with no necessary
connection to the arithmetic of elliptic curves
of the same conductor.
This demystifies the CM numerology that pervaded
earlier analyses~\cite{vquad-exclusion}.
In particular, PSLQ searches against CM period
bases at discriminant~$-11$ return null
\emph{structurally}, not merely at the current
precision.
\end{remark}

\subsection{Singularity structure}

The finite singularities of~\eqref{eq:ode}
are the roots of $a(x)$:
\begin{equation}\label{eq:sing}
  s_{1,2} = \frac{-1 \pm i\sqrt{11}}{6},
  \qquad \mathrm{disc}(a)
  = 1 - 4 \cdot 3 \cdot 1 = -11.
\end{equation}
Both are regular singular points. The point
$x=\infty$ is an irregular singular point of
Poincar\'e rank~1 (the coefficient $c(x)/a(x)$
has a pole of order~2 at $\infty$).

We define $\Vq = M_{11}$ as the
$(1,1)$-entry of the connection matrix
expressing the canonical recessive solution
at $+\infty$ in terms of the analytic
Frobenius basis at $s_1$; this is the
natural normalization for the Stokes
connection coefficient.

The WKB analysis at $\infty$ gives
$y \sim x^{\mu_\pm} e^{\sigma_\pm x}$ with
\begin{equation}\label{eq:wkb}
  \sigma_\pm = \pm \frac{1}{\sqrt{3}},
  \qquad
  \mu_\pm = -1 \mp \frac{1}{6\sqrt{3}}.
\end{equation}
The Stokes gap is
$\sigma_+ - \sigma_- = 2/\sqrt{3}$.


% ══════════════════════════════════════════════════════════════════
%  3. APPARENT SINGULARITY EXCLUSION
% ══════════════════════════════════════════════════════════════════
\section{Apparent singularity exclusion}
\label{sec:frobenius}

\begin{theorem}[Apparent singularity exclusion
for $\Vq$]
\label{thm:exclusion2}
Let $\Vq$ be the Stokes connection
coefficient $M_{11}$ of the ODE~\eqref{eq:ode}.
\begin{enumerate}
\item[\textup{(i)}]
  At both finite singularities
  $s_{1,2}=(-1\pm i\sqrt{11})/6$,
  the indicial equation is $\rho^2=0$,
  giving a double root $\rho=0$.
\item[\textup{(ii)}]
  The monodromy around each $s_k$ fixes
  $M_{11}=\Vq$; the effective monodromy
  on the connection coefficient is trivial.
  In particular, $\Vq$ cannot be identified
  via CM elliptic periods, gamma values, or
  classical hypergeometric special values by
  any monodromy-based argument.
\end{enumerate}
\end{theorem}

\begin{proof}
\noindent\textbf{Step 1: Indicial exponents.}
Write $P(x) = a'(x)/a(x)$ and
$Q(x) = c(x)/a(x)$.
At a simple root $s_k$ of $a$:
\[
  p_0 = \lim_{x\to s_k}
  (x-s_k)\,P(x)
  = \lim_{x\to s_k}
  \frac{(x-s_k)\,a'(x)}{a(x)}
  = \frac{a'(s_k)}{a'(s_k)}
  = 1,
\]
since $a(x) \sim a'(s_k)(x-s_k)$
near a simple root. And
\[
  q_0 = \lim_{x\to s_k}
  (x-s_k)^2\,Q(x)
  = \lim_{x\to s_k}
  \frac{(x-s_k)^2\,c(x)}{a(x)}
  = 0,
\]
since $a(x)$ has only a simple zero
at $s_k$, making $Q(x)$ have only a
simple pole and $(x-s_k)^2 Q(x) \to 0$.

The indicial equation is therefore
\[
  I(\rho) = \rho(\rho-1)
  + p_0\rho + q_0
  = \rho^2 = 0,
\]
giving a double root $\rho = 0$
at both $s_1$ and $s_2$.
This establishes part~(i).

\medskip
\noindent\textbf{Step 2: Analytic
first solution.}
For $\rho = 0$, the Frobenius ansatz
$y_1 = \sum_{n=0}^\infty c_n \zeta^n$
(with $\zeta = x - s_k$, $c_0 = 1$)
satisfies the recurrence
\[
  n^2 c_n = -\sum_{j=0}^{n-1}
  \bigl[(j+1)p_{n-1-j}
  + q_{n-1-j}\bigr] c_j,
  \quad n \geq 1,
\]
where $p_j$, $q_j$ are the Taylor
coefficients of $P(x)(x-s_k)$ and
$Q(x)(x-s_k)^2$ at $\zeta=0$.
Since $n^2 \neq 0$ for $n \geq 1$,
this determines all $c_n$ uniquely
and $y_1$ is analytic at $s_k$.

\medskip
\noindent\textbf{Step 3: Exact ODE form
and Wronskian.}
The identity $b(x) \equiv a'(x)$ gives
$a(x)\,y'' + a'(x)\,y'
= \frac{d}{dx}[a(x)\,y']$,
so the ODE~\eqref{eq:ode} is equivalent
to the exact form~\eqref{eq:exact-ode}.
By Abel's theorem, the Wronskian of any
two solutions satisfies
\[
  W[y_1, y_2](x) = \frac{C}{a(x)}
\]
for a constant $C$.
Near $s_k$, $a(x) \sim a'(s_k)\zeta$,
so $W \sim C/(a'(s_k)\zeta)$.

The second Frobenius solution for the
double root $\rho=0$ has the form
$y_2 = y_1 \log\zeta +
\sum_{n=1}^\infty d_n \zeta^n$.
One computes
\[
  W[y_1, y_2]
  = y_1^2/\zeta + O(1),
\]
consistent with the Abel formula.
The logarithm appears only in $y_2$,
not in $y_1$.

\medskip
\noindent\textbf{Step 4: Trivial
monodromy on $M_{11}$.}
The monodromy around $s_k$ acts as
\[
  \begin{pmatrix} y_1 \\ y_2
  \end{pmatrix}
  \mapsto
  \begin{pmatrix} 1 & 2\pi i \\
  0 & 1
  \end{pmatrix}
  \begin{pmatrix} y_1 \\ y_2
  \end{pmatrix}.
\]
The connection coefficient
$M_{11} = \Vq$ is the coefficient of
the analytic solution $y_1$ in the
expansion of the recessive solution at
$\infty$. Since monodromy around $s_k$
fixes $y_1$, the coefficient $M_{11}$
is unchanged. The effective monodromy on
$\Vq$ is therefore trivial,
establishing part~(ii).

\medskip
\noindent\textbf{Numerical certificate.}
The identity $b(x) \equiv a'(x)$
was verified symbolically.
The resulting values $p_0 = 1$ and
$q_0 = 0$ at both $s_{1,2}$ were
confirmed numerically at
\texttt{dps}$=150$:
\begin{align*}
  p_0(s_1) &= 1\;\text{ exactly},\quad
  q_0(s_1) = 0\;\text{ exactly},\\
  p_0(s_2) &= 1\;\text{ exactly},\quad
  q_0(s_2) = 0\;\text{ exactly}.
\end{align*}
Verification script:
\texttt{verify\_frobenius\_apparent.py};
results in\\
\texttt{frobenius\_apparent\_verification.json}.
\end{proof}

\begin{remark}
The content of Theorem~\ref{thm:exclusion2}
is not a computational exclusion but a
\emph{structural} one. No increase in PSLQ
precision can overcome the obstruction: the
monodromy at the finite singularities is
exactly trivial by algebra, and the CM
discriminant $-11$ is simply the discriminant
of the characteristic polynomial
$a(x)=3x^2+x+1$.
\end{remark}


% ══════════════════════════════════════════════════════════════════
%  4. PAINLEVÉ CLASSIFICATION
% ══════════════════════════════════════════════════════════════════
\section{Painlev\'e classification}
\label{sec:painleve}

The ODE~\eqref{eq:ode} has two regular
singularities (at $s_1,s_2$, both apparent)
and one irregular singularity of rank~1
(at $x=\infty$). This singularity
configuration on $\mathbb{P}^1$ is precisely
the setting of the Jimbo--Miwa--Ueno
isomonodromic deformation
theory~\cite{jimbo1981}.

\begin{theorem}[Painlev\'e~V ($D_5^{(1)}$) classification]
\label{thm:painleve}
The isomonodromic deformation family
containing the ODE~\eqref{eq:ode} is
governed by the fifth Painlev\'e equation
with parameters
\[
  \alpha = \frac{1}{6},\quad
  \beta = 0,\quad
  \gamma = 0,\quad
  \delta = -\frac{1}{2},
\]
so that $\beta=\gamma=0$ while
$\delta=-\tfrac12\neq0$. Its space of
initial conditions is the Sakai rational
surface of type $D_5^{(1)}$, with affine
Weyl symmetry
$W(A_3^{(1)})\cong\widetilde{\mathfrak{S}}_4$.
The vanishing $\beta = \gamma = 0$ reflects
the fact that both finite singularities are
apparent; $z=0$ is an ordinary point with
exponents $\{0,1\}$. Since $\delta\neq0$,
the equation lies off the PV$\to$PIII
confluence and does not reduce to
PIII($D_6$); the absence of any $z^{1/2}$
ramification excludes PIII($D_7$) as well.
\end{theorem}

\begin{proof}
The singularity data is: two regular
singularities with formal monodromy
exponents $\theta_1 = \theta_2 = 0$
(apparent), and one irregular singularity
at $\infty$ with Stokes data
$(\sigma_+, \sigma_-, \mu_+, \mu_-)$
given by~\eqref{eq:wkb}.

Following the Jimbo--Miwa--Ueno
classification~\cite{jimbo1981}, the
configuration
$\{2 \text{ regular (apparent)} +
1 \text{ irregular (rank 1)}\}$ on
$\mathbb{P}^1$ is the isomonodromy setting
of the fifth Painlev\'e equation, whose
space of initial conditions is the Sakai
surface $D_5^{(1)}$ with symmetry
$W(A_3^{(1)})$.
The parameters are determined by the
formal monodromy at $\infty$:
\[
  \alpha = \tfrac{1}{2}(\mu_+ - \mu_-)
  = \tfrac{1}{2} \cdot \tfrac{1}{3\sqrt{3}}
  \cdot \sqrt{3}
  = \tfrac{1}{6},
  \quad \delta = -\tfrac{1}{2}\sigma_+^2
  = -\tfrac{1}{2}.
\]
Since $\theta_1 = \theta_2 = 0$,
$\beta = \gamma = 0$.
Finally, the PV$\to$PIII degeneration of
DLMF~\S32.2(vi), Eq.~32.2.33, is a scaling
limit in which the PV $\delta$-term tends to
zero; since $\delta=-\tfrac12$ is fixed and
nonzero, no such limit is taken and the
equation does not reduce to PIII($D_6$).
Equivalently, $z=0$ is an ordinary point
(exponents $\{0,1\}$) with no $z^{1/2}$
branching, which excludes the ramified
PIII($D_7$).
\end{proof}

\begin{remark}[Connection coefficient]
The Jimbo connection formula expresses the
Stokes connection coefficient $M_{11}$ in
terms of the monodromy data. For generic
parameters, this involves products of gamma
functions evaluated at the formal monodromy
exponents. In the present case
$\theta_1 = \theta_2 = 0$ (apparent
singularities), the standard connection
formula degenerates, and the resulting
expression for $M_{11}$ is not immediately
computable in closed form. This is
consistent with the failure of all tested
PSLQ identification bases.
\end{remark}

\begin{remark}[Why not Painlev\'e~III]
The equation arises as a parameter
specialization ($\beta=\gamma=0$) of PV but
does \emph{not} further reduce to PIII. The
PV$\to$PIII confluence (DLMF~\S32.2(vi),
Eq.~32.2.33) is a scaling limit sending the
PV $\delta$-term to zero; here
$\delta=-\tfrac12\neq0$ and no such limit is
taken, so the equation remains genuinely PV
on $D_5^{(1)}$. The ramified PIII($D_7$)
reading is independently excluded because
$z=0$ is ordinary (integer exponents
$\{0,1\}$) with no $z^{1/2}$ branching.
\end{remark}

\begin{remark}[Generic, non-classical orbit]
The formal monodromy at infinity
$\theta_\infty=\sigma_+-\sigma_-=2/\sqrt{3}$
is irrational, so the parameters do not
satisfy the integer/resonance conditions
that characterize the Riccati, algebraic,
and classical special-function solution
families of PV under the $W(A_3^{(1)})$
action. Together with the exhaustive
PSLQ-null evidence of \S\ref{sec:stokes},
this indicates that $\Vq$ is a
\emph{generic, non-classical} PV
transcendent; we do not prove this and
state it as an expectation
(cf.\ Conjecture~\ref{conj:S_id}).
\end{remark}


% ══════════════════════════════════════════════════════════════════
%  5. RESURGENCE ANALYSIS
% ══════════════════════════════════════════════════════════════════
\section{Resurgence analysis}
\label{sec:resurgence}

\subsection{Formal asymptotic series}

The recessive solution at $x=+\infty$ has the
formal asymptotic expansion
\begin{equation}\label{eq:formal}
  y_{\mathrm{rec}}(x) \sim
  x^{\mu_-} e^{-x/\sqrt{3}}
  \sum_{n=0}^{\infty} a_n\,x^{-n},
  \qquad a_0 = 1,
\end{equation}
where $\mu_- = -1 + 1/(6\sqrt{3})$.
The coefficients $a_n$ are determined by
substituting~\eqref{eq:formal}
into~\eqref{eq:ode} and matching powers of
$x^{-n}$; this produces a two-term recurrence
with factorial growth $|a_n| \sim C \cdot n!
\cdot \xi_0^{-n}$ for a constant $C$,
confirming that the series is Gevrey-$1$
divergent.

\subsection{Borel transform}

The Borel transform
$\hat{y}(\xi) = \sum_{n=0}^{\infty}
a_n\,\xi^n / n!$
has radius of convergence
$|\xi_0| \approx 2/\sqrt{3} \approx 1.15470$,
the Stokes gap.
This value nearly coincides with the Stokes gap
$\sigma_+ - \sigma_- = 2/\sqrt{3}$ from
\eqref{eq:wkb}, as expected from the general
relation between exponential scales and Borel
singularities in resurgent systems.

\begin{theorem}[Borel singularity]
\label{thm:borel}
The Borel transform $\hat{y}(\xi)$ of
the formal series~\eqref{eq:formal} has
a branch-point singularity at
$\xi_0 \approx 2/\sqrt{3}$ (numerically
$\xi_0 = 1.15472\ldots$, a near-miss at the
$2\times10^{-5}$ level; we do not claim an exact
algebraic value, consistent with the generic,
non-classical character of the transcendent)
on the negative real Borel axis.
\begin{enumerate}
\item[\textup{(i)}]
  The Domb--Sykes ratio
  $r_n = |a_n/a_{n-1}|$ satisfies
  $r_n \to 1/\xi_0 \approx \sqrt{3}/2$,
  confirmed by Richardson extrapolation
  to five significant figures.
\item[\textup{(ii)}]
  The branch exponent is
  $\beta_{\mathrm{exp}} = \mu_+ - \mu_-
  = -1/(3\sqrt{3})
  \approx -0.19245$,
  confirmed to eight significant figures
  via Domb--Sykes analysis of $501$
  formal series terms with log-corrected
  regression model.
\item[\textup{(iii)}]
  The singularity structure
  (logarithmic branch point, not a simple
  pole) is consistent with the generic
  resurgence structure of Painlev\'e~V
  transcendents.
\end{enumerate}
\end{theorem}

\begin{proof}
(i) We compute $501$ coefficients $a_n$ of
the formal series via the Riccati recurrence
at \texttt{dps}$=200$, then form the ratio
sequence $r_n = |a_n/a_{n-1}|$.
Richardson extrapolation of order~3 applied
to $1/r_n$ gives
$\xi_0 \approx 1.15472$, matching
$2/\sqrt{3} = 1.15470\ldots$ to five
significant figures.

(ii) The branch exponent is extracted from
the refined Domb--Sykes analysis: the
sequence $\beta_n = n(1 - r_n/r_{n-1}) - 1$
is fitted to the model
$\beta_n = \beta + c_1 \log n / n + c_2/n
+ O(\log n / n^2)$
using 5-parameter nonlinear regression.
The result
$\beta_{\mathrm{est}} = -0.192450088$
agrees with
$-1/(3\sqrt{3}) = -0.192450090\ldots$
to eight significant figures.

(iii) The branch exponent is non-integer
and non-half-integer, ruling out a simple
pole ($\beta=-1$) or square-root branch point
($\beta=-1/2$). The logarithmic branch-point
structure is the generic singularity type
for Borel transforms of Painlev\'e~V
trans-series, where the exponent equals
the difference of formal monodromy
exponents at $\infty$.
\end{proof}

A rigorous proof of Borel summability for this class of
continued-fraction formal series, in the sense of
Costin~\cite{costin1998}, would require establishing
Gevrey-$1$ bounds on the coefficient sequence~$a_n$;
we leave this as an open problem and treat the present
result as strong numerical evidence.


% ══════════════════════════════════════════════════════════════════
%  6. THE STOKES CONSTANT
% ══════════════════════════════════════════════════════════════════
\section{The Stokes constant}
\label{sec:stokes}

The Stokes constant $S$ governs the
exponentially small discontinuity across
the Stokes line:
\[
  y_{\mathrm{rec}}^{(\text{Stokes}+)}
  - y_{\mathrm{rec}}^{(\text{Stokes}-)}
  = S \cdot y_{\mathrm{dom}}.
\]
By the Dingle late-term
formula~\cite{dingle1973},
\begin{equation}\label{eq:dingle}
  S_n = \frac{a_n \cdot \Gamma(\beta_{\mathrm{exp}})
  \cdot \xi_0^{n+\beta_{\mathrm{exp}}}}
  {(-1)^n \cdot \Gamma(n+\beta_{\mathrm{exp}})}
  \;\longrightarrow\; S
  \quad\text{as } n\to\infty.
\end{equation}

\begin{proposition}[Stokes constant estimate]
\label{prop:stokes}
Using $501$ formal series terms with
$\beta_{\mathrm{exp}} = -1/(3\sqrt{3})$ and
$\xi_0 \approx 2/\sqrt{3}$:
\begin{enumerate}
\item[\textup{(i)}]
  Dingle--Richardson${}^2$ extraction in
  $1/(n+\beta)$ gives
  $S \approx 0.43770528$.
\item[\textup{(ii)}]
  Split-half consistency (first 250 vs.\
  last 251 terms) confirms $8$ significant
  digits.
\item[\textup{(iii)}]
  PSLQ searches against $14$ independent
  bases (including $\Gamma(k/6)$ products,
  $\pi\sqrt{3}$ combinations, $\exp(\pi/\sqrt{3})$,
  and algebraic numbers over
  $\Q(\sqrt{3},\sqrt{11})$) all return null.
\end{enumerate}
The quoted $8$-digit stability is based on
split-half consistency and the observed plateau
of the Richardson-extrapolated sequence; we do
not claim a rigorous error bound.
\end{proposition}

The bridge equation relating the Stokes constant~$S$ to
alien derivatives of the formal solution, in the sense of
\'Ecalle's resurgence theory~\cite{ecalle1981}, is not
established here; Conjecture~\ref{conj:S_id} provides
a conjectural identification via the Jimbo~\cite{Jimbo1982}
formula.

\begin{remark}
The $8$-digit precision of $S$ is insufficient
for definitive PSLQ identification, which
typically requires $15+$ digits for a
reliable hit. Reaching this precision would
require $2000+$ formal series terms at
\texttt{dps}$=500$, or fundamentally
different methods such as conformal
Borel--Pad\'e resummation or Berry--Howls
hyperasymptotics. The identification of $S$
in closed form is an open problem.
\end{remark}


% ══════════════════════════════════════════════════════════════════
%  7. DISCUSSION AND OPEN PROBLEMS
% ══════════════════════════════════════════════════════════════════
\section{Discussion and open problems}
\label{sec:discussion}

\subsection{Summary of the identification
obstruction}

The structural exclusion of
Theorem~\ref{thm:exclusion2} and the
Painlev\'e classification of
Theorem~\ref{thm:painleve} together
provide a complete explanation for the
failure of identification searches
for $\Vq$:
\begin{enumerate}
\item The finite singularities are apparent
  (trivial monodromy), so $\Vq$ cannot arise
  as a period or CM value at these points.
\item The irregular singularity at $\infty$
  places the problem in the confluent-Heun /
  Painlev\'e~V ($D_5^{(1)}$) sector, outside
  the Gauss hypergeometric framework.
\item As a generic, non-classical PV
  transcendent (off the Riccati, algebraic,
  and classical special-function orbits),
  its connection coefficient is not expected
  to admit a $\Gamma$/Barnes closed form.
\item The Jimbo formula search (iteration~23)
  established that the formal WKB exponent
  $\sigma_\pm = \pm 1/\sqrt{3}$ is distinct
  from the connection parameter
  $\sigma_{\mathrm{conn}}$ required by the
  Jimbo formula, precisely locating the
  obstruction to closed-form identification
  of~$S$.
\end{enumerate}

\subsection{The connection parameter and
Jimbo formula}

\begin{conjecture}[Identification of $S$]
\label{conj:S_id}
The Stokes constant $S\approx
0.43770528\ldots$ is given by the
Jimbo (1982) \cite{Jimbo1982} formula
\[
  S = 2i\sin(\pi\sigma_{\mathrm{conn}})
  \,\frac{\Gamma(1-\sigma_{\mathrm{conn}})^2}
  {\Gamma(1+\sigma_{\mathrm{conn}})^2},
\]
where $\sigma_{\mathrm{conn}}$ satisfies
$\mathrm{tr}(M_1 M_2) =
2\cos(2\pi\sigma_{\mathrm{conn}})$
for the monodromy matrices of the
associated Riemann--Hilbert problem,
and is a transcendental function of the
accessory parameter
\[
  q = \frac{5+i\sqrt{11}}{54}.
\]
Numerically,
\[
  \sigma_{\mathrm{conn}} =
  0.06087689082546\ldots
\]
($15$ digits, obtained by inverting the
Jimbo formula from the $8$-digit Stokes
constant $S$; limited by $S$ precision).
An exhaustive PSLQ search over $12$
bases including algebraic extensions of
$\mathbb{Q}(\sqrt{11})$, arctan/$\pi$
combinations, CM gamma values, Dedekind
eta quotients at $\tau = (-1+\sqrt{-11})/2$,
and the accessory parameter $q$ itself
found no identification at $8$-digit
precision.

An exhaustive search over $245$ formula
variants (nine theta-parameter
identifications, twenty formula types,
thirty-five theta-free formulas) found
no match to $S$ within $10^{-4}$,
confirming that the WKB exponent
$\sigma_\pm = \pm 1/\sqrt{3}$ is
distinct from $\sigma_{\mathrm{conn}}$.

Identification of $\sigma_{\mathrm{conn}}$
likely requires the $\tau$-function of
the Painlev\'e~V equation ($D_5^{(1)}$) at the
accessory parameter $q$; in the CM case
this may reduce to classical theta
series via the conformal block approach
of Gamayun--Iorgov--Lisovyy
\cite{GIL2012}.
\end{conjecture}

\begin{remark}[CFT approach]
The $\tau$-function approach of
\cite{GIL2012} expresses the Painlev\'e
connection parameter as the zero of a
Fourier series with Nekrasov partition
function coefficients. In the CM case
(discriminant $-11$), these may reduce
to theta series of the CM elliptic curve
with $j$-invariant $-32768$,
potentially yielding a closed form for
$\sigma_{\mathrm{conn}}$ and hence $S$.
\end{remark}

\subsection{Further open problems}

\begin{enumerate}
\item \textbf{Transcendence.}
  Is $\Vq$ transcendental? The structural
  exclusion from hypergeometric families and
  the non-D-finite evidence
  of~\cite{vquad-exclusion} suggest
  transcendence, but a proof would require
  fundamentally new methods.

\item \textbf{Universality of the $b=a'$
  structure.}
  Does the $b(x)=a'(x)$ structure persist
  when $a(n)$ is a non-constant polynomial?
  For degree-$d$ PCFs with polynomial $a(n)$,
  the Wallis characteristic polynomial has
  degree $2d$ and the apparent singularity
  condition imposes $2d$ constraints on the
  ODE coefficients---characterizing when
  these are satisfied is the general question.
\end{enumerate}


% ══════════════════════════════════════════════════════════════════
%  ACKNOWLEDGEMENTS
% ══════════════════════════════════════════════════════════════════
\section*{Acknowledgements}

Computations were performed using
\texttt{mpmath}~\cite{mpmath} for
arbitrary-precision arithmetic.
Supplementary scripts and data:
\href{https://github.com/papanokechi/pcf-research}{%
github.com/papanokechi/pcf-research}.

\smallskip
\noindent\textit{AI assistance disclosure.}
The preparation of this manuscript involved
assistance from Claude
(\texttt{claude-sonnet-4-6}, Anthropic) and
GitHub Copilot (Microsoft).
All mathematical results were conceived,
directed, and validated by the author.


% ══════════════════════════════════════════════════════════════════
%  REFERENCES
% ══════════════════════════════════════════════════════════════════
\begin{thebibliography}{99}

\bibitem{apery1979}
R.~Ap\'ery,
\emph{Irrationalit\'e de $\zeta(2)$ et
$\zeta(3)$},
Ast\'erisque \textbf{61} (1979), 11--13.

\bibitem{dingle1973}
R.~B.~Dingle,
\emph{Asymptotic Expansions: Their
Derivation and Interpretation},
Academic Press, London, 1973.

\bibitem{GIL2012}
O.~Gamayun, N.~Iorgov, and O.~Lisovyy,
\emph{Conformal field theory of
Painlev\'e~VI},
J.~High Energy Phys.\ \textbf{2012}
(2012), no.~10, 038.
arXiv:1207.0787.

\bibitem{Jimbo1982}
M.~Jimbo,
\emph{Monodromy problem and the boundary
condition for some Painlev\'e equations},
Publ.\ Res.\ Inst.\ Math.\ Sci.\
\textbf{18} (1982), no.~3, 1137--1161.

\bibitem{jimbo1981}
M.~Jimbo, T.~Miwa, and K.~Ueno,
\emph{Monodromy preserving deformation of
linear ordinary differential equations with
rational coefficients. I.},
Physica D \textbf{2} (1981), 306--352.

\bibitem{lorentzen2008}
L.~Lorentzen and H.~Waadeland,
\emph{Continued Fractions: Convergence Theory},
2nd~ed., Atlantis Studies in Mathematics,
vol.~1, Atlantis Press, 2008.

\bibitem{mpmath}
F.~Johansson et~al.,
\texttt{mpmath}: a Python library for
arbitrary-precision floating-point arithmetic
(version 1.3), 2023.
\url{http://mpmath.org/}

\bibitem{raayoni2021}
G.~Raayoni et~al.,
\emph{Generating conjectures on fundamental
constants with the Ramanujan Machine},
Nature \textbf{590} (2021), 67--73.

\bibitem{vquad-exclusion}
papanokechi,
\emph{Spectral classes of polynomial
continued fractions: Arithmetic
classification and constant identification},
submitted, 2026.
Manuscript ID 268704746.

\bibitem{ronveaux1995}
A.~Ronveaux (ed.),
\emph{Heun's Differential Equations},
Oxford University Press, 1995.

\bibitem{okamoto1987}
K.~Okamoto,
Studies on the Painlev\'e equations~I,
\emph{Ann.\ Mat.\ Pura Appl.}\ \textbf{146} (1987), 337--381.

\bibitem{costin1998}
O.~Costin,
On Borel summation and Stokes phenomena for rank-1 nonlinear
systems of ordinary differential equations,
\emph{Duke Math.\ J.}\ \textbf{93} (1998), 289--344.

\bibitem{ecalle1981}
J.~\'Ecalle,
\emph{Les fonctions r\'esurgentes}, Vols.~I--III,
Publ.\ Math.\ d'Orsay, 1981--1985.

\end{thebibliography}

\end{document}
